% Canonical Paper IV section.
A quotient forgets distinctions by definition.  It becomes a computationally
meaningful state only after a future action and an observation have been
specified.  This section gives the exact group-theoretic test for the
projective quotient $G/H$.

\subsection{Equivariant quotients and future sufficiency}

Let a monoid $M$ act on a set $X$ from the left, let $q:X\to Q$ be a
surjection, and suppose an action of $M$ on $Q$ satisfies
\[
  q(mx)=m q(x).
\]
Let $o:Q\to Y$ be an observation.  For $x,x'\in X$, define
\[
  x\equiv_{M,o}x'
  \quad\Longleftrightarrow\quad
  o(mq(x))=o(mq(x'))\quad\text{for every }m\in M.
\]

\begin{theorem}[Equivariant quotient as a residual state]
\label{thm:equivariant-quotient-residual}
If $q(x)=q(x')$, then $x\equiv_{M,o}x'$.  If the identity belongs to $M$ and
$o$ is injective, then the converse also holds.  Under those hypotheses the
future-residual quotient of $X$ is exactly $Q$.
\end{theorem}

\begin{proof}
Equivariance gives $q(mx)=mq(x)=mq(x')=q(mx')$, and applying $o$ proves the
first statement.  For the converse, test the identity continuation.  Then
$o(q(x))=o(q(x'))$, and injectivity of $o$ gives $q(x)=q(x')$.
\end{proof}

The theorem does not say that every quotient is sufficient.  It says that
sufficiency follows from a declared descended action; minimality additionally
uses an observation that can distinguish quotient states.

\subsection{Left continuation on the bivaluation quotient}

For the chronology convention of Equation~\eqref{eq:chronological-evaluation},
a future operator $k$ acts on the current operator $g$ by left multiplication:
$g\mapsto kg$.  The right-coset quotient therefore has the canonical action
\[
  k\cdot(gH)=(kg)H.
\]

\begin{corollary}[Exact left-continuation state]
\label{cor:left-continuation-exact}
Let all elements of $G$ be allowed as left futures, include the identity
future, and observe the resulting point of $G/H$ through an injective map.
Then two operators have the same future behavior exactly when they lie in the
same right coset:
\[
  g\equiv g'\quad\Longleftrightarrow\quad gH=g'H.
\]
Thus $G/H$ is the canonical minimal exact state space for this continuation
experiment.
\end{corollary}

In particular, the $H$-coordinate discarded by projective condensation is
invisible to every left continuation whose observation factors through the
same quotient.

\subsection{Right continuation and the normalizer boundary}

Opposite chronology replaces $g$ by $gk$.  Unlike left multiplication, this
does not automatically act on right cosets.

\begin{theorem}[Descent of one right update]
\label{thm:right-update-descent}
For a fixed $k\in G$, the formula
\[
  R_k(gH)=gkH
\]
defines a well-defined map $G/H\to G/H$ if and only if
\begin{equation}
  k^{-1}Hk\subseteq H.
  \label{eq:right-descent-inclusion}
\end{equation}
It is a bijection with inverse induced by $k^{-1}$ if and only if
$k\in N_G(H)$.
\end{theorem}

\begin{proof}
If $gH=ghH$ with $h\in H$, well-definedness requires
$ghkH=gkH$.  Cancelling $gk$ on the left gives
$k^{-1}hk\in H$ for every $h\in H$, which is
\eqref{eq:right-descent-inclusion}.  The same argument for $k^{-1}$ gives the
reverse inclusion.  Both hold exactly when $k^{-1}Hk=H$, equivalently
$k\in N_G(H)$.
\end{proof}

\begin{corollary}[Reversible right continuation]
\label{cor:right-group-normalizer}
A subgroup $L\le G$ acts on $G/H$ by right multiplication if and only if
$L\subseteq N_G(H)$.
\end{corollary}

For an infinite group, the inclusion in
\eqref{eq:right-descent-inclusion} may be strict, so the single-update and
reversible statements must not be conflated.  For finite $H$, equal
cardinality turns the inclusion into equality.

\begin{corollary}[Chronology-dependent visibility]
Suppose $h\in H$ is discarded by $g\mapsto gH$.  Under left continuation one
always has $kghH=k gH$, so $h$ remains invisible.  Under a right continuation
$k$ for which $k^{-1}hk\notin H$, the two representatives $g$ and $gh$ no
longer determine the same quotient state after continuation.  The discarded
frame direction has become observable.
\end{corollary}

This is the first rigorous bridge between projective process residue and a
future-sensitive residual.  The bridge is not the cardinality of $H$ alone;
it is the interaction between $H$ and the declared continuation side.

\subsection{The split-torus normalizer in rank one}

Take the reference pair $(0,\infty)$.  Then
\[
  H=\{z\mapsto az:a\in K^\times\}.
\]
Let $J(z)=1/z$; the sign is irrelevant projectively.

\begin{proposition}[Normalizer of the ordered-pair stabilizer]
\label{prop:split-torus-normalizer}
If $|K|>2$, then
\[
  N_G(H)=H\sqcup JH.
\]
Equivalently, a projective map normalizes $H$ exactly when it preserves the
unordered set $\{0,\infty\}$; it either fixes the two points or interchanges
them.
\end{proposition}

\begin{proof}
Because $|K|>2$, choose $a\ne1$.  The map $z\mapsto az$ has fixed set exactly
$\{0,\infty\}$.  Hence the common fixed set of $H$ is $\{0,\infty\}$.  A
normalizer sends this common fixed set to itself.  If it fixes both points it
lies in $H$; if it swaps them, composition with $J$ fixes both and lies in
$H$.  Conversely, $H$ normalizes itself and $J(z)=1/z$ conjugates
$z\mapsto az$ to $z\mapsto a^{-1}z$.
\end{proof}

Thus most right projective futures do not descend to the bivaluation quotient.
Remembering only the rank-one projector is sufficient for the left chronology,
but generally insufficient for arbitrary two-sided process composition.

\subsection{Sufficient, exact, and minimal condensations}

\begin{definition}[Future-sufficient condensation]
For a past space $P$, a continuation class $M$, and an observation $\Obs$, a
map $q:P\to Q$ is future-sufficient if there exist descended transitions on
$Q$ and an output map $o:Q\to Y$ such that every observation after every
continuation can be computed from $q(p)$.
\end{definition}

\begin{definition}[Exact and minimal condensation]
A future-sufficient condensation is exact when
$q(p)=q(p')$ precisely for pasts indistinguishable by all allowed futures.  It
is minimal when every other exact deterministic state representation maps
surjectively onto it.
\end{definition}

Theorem~\ref{thm:equivariant-quotient-residual} proves exactness for the
projective left-action experiment.  Section~\ref{sec:contextual-residuals}
gives the general residual construction and proves minimality without assuming
a group quotient.

\subsection{An exact-transport obstruction}

A tempting way to label a path of states is to choose frames $h_t\in G$ and
put
\[
  r_t=h_t h_{t-1}^{-1}.
\]
This transport contains no independent edge information.

\begin{proposition}[Exact edge labels telescope]
\label{prop:exact-transport-telescope}
For $s<t$,
\[
  r_t r_{t-1}\cdots r_{s+1}=h_t h_s^{-1}.
\]
Every closed lifted path therefore has trivial product.
\end{proposition}

\begin{proof}
All intermediate factors cancel pairwise.
\end{proof}

A nontrivial projective process holonomy consequently requires more than an
endpoint potential.  Possible additional data include an independent
stabilizer-valued factor in each edge, a connection on $G\to G/H$, a
multivalued frame, or rewrite $2$-cells.  None is supplied automatically by the
static quotient.
